\section{逃逸是相对运动的边界条件}
\subsection{无穷远剩余速度与表面速度}

在固定球形天体外，单位质量能量 $\varepsilon=v^2/2-\mu/r$。若轨道可向外延伸到无穷远，则 $\varepsilon=v_\infty^2/2$，所以局部速率满足 $v^2=v_\infty^2+2\mu/r$。$v_\infty$ 称双曲余速；逃逸的临界值是 $v_\infty=0$，对应当地逃逸速度 $v_{\rm esc}=\sqrt{2\mu/r}$。

能量非负只说明轨道存在延伸到无穷远的分支。若当前速度径向向内且角动量为零，物体会先撞上中心，不能把“存在数学上的出射分支”当作实际逃逸。对有尺寸的天体，还要检查近地点是否在表面之外；对能反弹、能穿透或会被吸收的物体，后续边界也不同。

两体都能运动时，在质心系使用相对速度 $\bm u=\bm v_1-\bm v_2$，相对引力参数为 $G(m_1+m_2)$。两体分离到无穷远的能量阈值由相对运动决定，整体质心在实验室中平动得再快，也不会单独帮助两体摆脱相互束缚。

\subsection{相落时间的正、零、负能量分支}

纯径向向内的初速为 $u_0$、初间距 $l$，能量式给 $\dot r=-\sqrt{u_0^2+2\mu(1/r-1/l)}$。相遇时间应从此负根换成正的积分
\[t=\int_0^l\frac{\dd r}{\sqrt{u_0^2+2\mu(1/r-1/l)}}.\]
这个表达式对三种能量都成立。若令 $k=2\mu/l-u_0^2$，$k>0$ 用三角换元，$k=0$ 直接幂积分，$k<0$ 用双曲函数换元。只写负能量的反三角函数而不检验初速范围，会在高速入射时得到非实数时间。


\section{双曲线、散射与碰撞截面}
\subsection{几何参数的正负约定}

标准双曲线 $x^2/a_h^2-y^2/b_h^2=1$ 中 $a_h,b_h>0$，$c_h^2=a_h^2+b_h^2$、$e=c_h/a_h>1$。渐近线倾角满足 $\tan\theta=b_h/a_h$，焦点到渐近线的垂距恰为 $b_h$。这个垂距是无穷远入射的冲量参数，并非近地点距离。

引力双曲线能量为正：$\varepsilon=\mu/(2a_h)$，而角动量满足 $h^2=\mu a_h(e^2-1)=\mu b_h^2/a_h$。若沿用活力公式的有符号半长轴，则应取 $a=-a_h<0$，不能把正的几何 $a_h$ 直接塞进椭圆的负能量式。

\typeout{ASTROFIG:fig:at4:hyperbola}\fig{fig:at4:hyperbola}{入射渐近线与最近点属于不同几何量}{\begin{tikzpicture}[scale=.85]
\draw[guide](-4.6,-1.118)--(1,-1.118);\draw[world,domain=-2:2,samples=110]plot({(1.5-cosh(\x))/1.5+1.11803399*1.11803399*sinh(\x)/1.5},{-1.11803399*(1.5-cosh(\x))/1.5+1.11803399*sinh(\x)/1.5});\fill(0,0)circle(2pt)node[above left]{$F$};\draw[flow](-4.2,-1.118)--(-3,-1.118)node[midway,below]{$v_\infty$};\draw[dim](1,-1.118)--(1,0)node[midway,right]{$b$};\draw[guide](0,0)--(.333,-.373);\node[right]at(.5,-.25){$r_p$};
\end{tikzpicture}}
图\ref{fig:at4:hyperbola}中的 $b$ 垂直于入射渐近线，而 $r_p$ 从焦点指向最近点。两者之差反映了引力聚焦，图上不可把这两段长度混用。\subsection{散射角与最近距离}

无穷远入射速度 $v_\infty$、冲量参数 $b$，给 $h=bv_\infty$。由能量与离心率式，
\[e=\sqrt{1+\frac{b^2v_\infty^4}{\mu^2}},\quad
r_p=\frac{h^2/\mu}{1+e}=\frac\mu{v_\infty^2}(e-1).\]
吸引散射的速度转向角满足 $\sin(\chi/2)=1/e$，等价于 $\tan(\chi/2)=\mu/(bv_\infty^2)$。这里比较的是入射与出射速度，而不是两条朝外径矢之间的张角。远处速率不变，方向改变；若两端还保留同一中心势，能量守恒不等于速度矢量守恒。

对于半径 $R$ 的吸收体，刚好擦到表面的轨道有 $r_p=R$。消去近点速度或离心率，得到最大冲量参数 $b_{\max}^2=R^2(1+2\mu/(Rv_\infty^2))$。入射截面是 $\pi b_{\max}^2$，通常大于几何截面 $\pi R^2$。该公式假定稀薄粒子、无相互遮挡、无阻力以及命中后确实被吸收。


\section{碰撞改变轨道时分两段计算}
\subsection{短时动量守恒与随后机械能守恒}

陨石与卫星粘在一起，碰撞时间很短，可忽略这段时间中心引力的冲量，因此两物体的总动量近似守恒。非弹性碰撞一般损失机械能，不能跨碰撞全过程写卫星机械能守恒。求出碰后速度矢量后，再把合体看作新的试验物体，计算后续的 $\varepsilon,h,r_p$。

即使陨石沿径向撞入，卫星原有的切向动量仍在。碰后切向速度减小，但一般不为零，因此不能只比较碰后总能量与表面势能；必须给表面剩余切向动能留出能量。以到达表面时径向速度平方非负为条件，常比先求完整椭圆参数更直接。

\subsection{接触等号和时间方向}

若题目允许擦碰，边界是表面处径向速度为零；要求进入或撞击具有有限法向速度，则需要严格不等式。得到一个阈值表达式以后，还要检查它是否为正：如果无径向陨石速度时，单是切向减速就已使轨道进入表面，所求最低径向速度就是零，而不是对一个负数强行开根。


\section{移动行星的引力助推}
\subsection{先在行星系散射，再回到恒星系}

在局部掠近期间，近似行星的恒星系速度 $\bm V$ 不变，忽略恒星在行星附近的差分场。入、出行星引力影响区时，先作速度变换 $\bm u_{\rm in}=\bm v_{\rm in}-\bm V$。行星系双曲散射保持余速大小 $u$，只把方向旋转一个允许角 $\delta$，然后变回 $\bm v_{\rm out}=\bm V+\bm u_{\rm out}$。

在这个近似中，恒星系单位质量能量变化为 $\bm V\cdot(\bm u_{\rm out}-\bm u_{\rm in})$。它可以正也可以负，取决于从行星的哪一侧绕过。助推是与移动天体交换动量和能量，不是静止引力场无代价增加同一惯性系中的机械能。

\typeout{ASTROFIG:fig:at4:frames}\fig{fig:at4:frames}{恒星系速度是行星平动与相对速度的矢量和}{\begin{tikzpicture}[scale=.85]
\draw[flow](0,0)--(2.5,0)node[midway,below]{$\bm V$};\draw[flow](2.5,0)--(3.5,1.8)node[midway,right]{$\bm u$};\draw[force](0,0)--(3.5,1.8)node[midway,above left]{$\bm v$};
\end{tikzpicture}}
图\ref{fig:at4:frames}的速度三角形适用于入射和出射两端。只给速度大小而不给绕行一侧和冲量参数，通常不能得到唯一的出射矢量。\subsection{最大转角不能任意取反向}

若行星半径、大气或安全条件规定近行星点 $r_p\ge r_{\min}$，对给定余速 $u$，允许转角不超过
\[\delta_{\max}=2\arcsin\frac1{1+r_{\min}u^2/\mu_p}.\]
有限正的 $r_{\min},u$ 给 $\delta_{\max}<\pi$。把相对速度无条件反转180度，会夸大助推能力，有时对应穿入行星或极限冲量参数为零的不可行路径。

从地表逃离行星、沿恒星轨道转移、再掠近另一行星，是几个不同的保守飞行段。先用地球系求脱离余速，再变到日心系；到火星附近再转到火星系；离开火星后才用日心能量判断逃逸太阳。最小发射速度还依赖发射方向、会合相位和最小掠近半径，不能只把两个逃逸速度标量相加。


\section{散射中的时间、捕获与动量交换}
\subsection{双曲近点前后的时间}

  正能量轨道可用双曲偏近点参数 $H$ 表示，取近点为 $H=0$，$r=a_h(e\cosh H-1)$。面积速度不变给 $t-t_p=\sqrt{a_h^3/\mu}(e\sinh H-H)$。其导数因子 $e\cosh H-1>0$，所以同一有向分支上的时间随 $H$ 单调增加。

  初位置在入射支，参数应取负值；到近点的剩余时间取正。把余弦双曲函数反函数得到的正值直接当成初始 $H$，会把入射位置误认成出射位置。与椭圆不同，双曲线没有周期，也不能用某个“完整双曲线面积”除以面积速度求飞行时间。
  
\subsection{静止天体的保守引力不能自行捕获}

  无穷远以非零余速入射的物体，相对能量为正。若只有固定中心的保守引力，它在远处出射时仍有同样的正能量；单纯靠轨道弯曲不能把双曲线变成椭圆。捕获需要阻力、碰撞、点火或第三个天体交换能量，或者改变了所选系统边界。

  若在近点逆行点火，局部速率为 $\sqrt{v_\infty^2+2\mu/r_p}$。减到当地抛物线速率 $\sqrt{2\mu/r_p}$ 只是零能量边界，要形成真正椭圆还需稍多减速。取最小近点能减少所需速度增量，但同时受到实体半径、大气、潮汐和设备约束。
  
\subsection{吸积物质带来的动量变化}

  天体吸收静止尘埃时，合并前后动量关系为 $M\bm V=(M+\dd M)(\bm V+\dd\bm V)$，保留一阶得 $M\dd\bm V=-\bm V\dd M$。因此仅就吸收静止物质而言，天体会减速；其动能减少量并不全由平动速度变化项表达，还应包括质量增加项。

  引力聚焦只决定哪些远处粒子能命中，并不取消它们原来的动量条件。要计算长期轨道演化，必须把粒子原来的日心速度、吸收过程与天体受外力一并纳入，不能把“吸积率为正”直接等同于“天体获得推力加速”。
  

\section{分段二体近似与完整可行性}
\subsection{同一个无穷远在不同近似中的含义}

  掠近行星时常把行星附近看成二体系统，把足够远处的相对运动看成直线渐近线。这个“行星系无穷远”是局部二体轨道外推得到的余速，不表示在真实太阳系中可以一路忽略太阳，直到空间位置真的趋于无穷。离开局部区域后，应把渐近余速与行星公转速度相加，接回日心运动。

  检验局部行星主导，应比较太阳对探测器和行星的差分加速度，而非太阳对二者共同的加速度。若局部尺度为 $\ell$、行星日心距离为 $r_p$，太阳潮汐量级为 $\mu_s\ell/r_p^3$，行星引力为 $\mu_p/\ell^2$。前者远小于后者要求
  \[\ell\ll r_p(\mu_p/\mu_s)^{1/3}.\]
  这是量级条件，不在此承诺精确的影响区边界。遇到非常慢、非常远或长期绕行的掠近，太阳扰动和行星曲线运动都可能积累到不可忽略，局部固定速度散射模型应重新检查。
  
\typeout{ASTROFIG:fig:at4:patch}\fig{fig:at4:patch}{局部行星散射与日心飞行在共同边界上接续}{\begin{tikzpicture}[scale=.85]
\draw[guide](0,0)circle(1.6);\fill(0,0)circle(4pt)node[below]{行星};\draw[flow](-3,-1)--(-1.4,-.8)node[midway,above]{入界};\draw[world](-1.4,-.8)..controls(0,-1)and(1.2,-.5)..(1.4,.8);\draw[flow](1.4,.8)--(2.8,1.9)node[midway,left]{出界};\node at(0,-2.2){边界上转换速度参考系};
\end{tikzpicture}}

  图\ref{fig:at4:patch}的虚线圈是近似区域示意，并非一个有力的实体边界。边界两侧的位置和惯性系速度必须连续；变化的是计算中保留哪一个主导天体。不能在入界时把相对速度直接替换成日心速度，或在出界时把行星速度忘记加回。

  条件不足的助推题可以先列出可确定量和剩余自由度。能量式确定余速大小，角动量或冲量参数确定转角，绕行一侧确定转角符号，行星公转速度则负责把结果变回日心系。缺少其中任一项，可能只能给速度范围或一个参数族。保留参数不是省略计算，而是防止凭空增加题设没有给出的条件。

  有限安全半径还会把可行速度范围缩小。对同一个余速，转角有上限；对同一个目标出射方向，还要问这个方向是否在允许转向扇区内。先在速度图中画出可转动的扇区，再最大化日心速度投影，比直接设定“末速度必与行星速度同向”更稳妥。若那个方向不在扇区内，最优值在扇区边界而不是理想平行方向。

  实际最值的结论应区分已取得的最小值、只可逼近的下确界，以及由必要条件得到的下界。抛物线捕获边界、无穷中间半径和极限180度反转分别说明了三种容易混淆的情形。最后回代检验表面、方向、相位和时间，是让代数结果成为物理解答的必要一步。
  
